% Introduction paragraph 3 (≈250w) — Meno-problem framing.
% For task t-mozv8oul (Biological Science) and BioSci's locked abstract/Intro-P3 slot.
% Anchors: Math Assumption (A4) / Lemma S1-exchangeability (bound), framework name REAL,
% motif set W, protection RareShield, per-seed score LRS(w).

\paragraph*{A Meno problem in empirical science.}
Plato's \emph{Meno} asks how one can search for what one does not know: if the object of inquiry is unknown, it cannot be recognised even if encountered. The preservation of previously unannotated rare cell subpopulations during single-cell batch integration instantiates this paradox in a concrete form. Every existing evaluation---scIB, scIB-E, RBET, CellANOVA---certifies preservation against an annotation. If a subpopulation has never been annotated, no score moves when it is destroyed, and no score moves when it is preserved; the object is simply not in the evaluator's vocabulary. Our resolution is to replace the \emph{semantic} criterion (does the population have a name?) with a \emph{syntactic} criterion (is there reproducible structure across independent observations?). Reproducibility is a label-free property, witnessed geometrically by local-density motifs $w \in \mathcal{W}$ that recur across batches and persist under admissible batch-effect deformations, and quantified by a per-motif survival score $\mathrm{LRS}(w)$. A motif whose cross-batch reproducibility predates integration but not its integrated image is diagnostic of erasure, regardless of whether any name has been attached to it. We validate the criterion on held-out known rare types (pancreatic $\epsilon$, pulmonary ionocytes, LAMP3$^+$ mregDC, TPEX, FOLR2$^+$ macrophages) by hiding their labels, confirming detection, and transferring the estimate to genuinely unannotated subpopulations under an exchangeability bound (Lemma~\ref{lem:exchangeability}). REAL detects; RareShield protects. The bound, not the field's habits, licenses the transfer.
